\subsection{碰撞}
\begin{tabularx}{\columnwidth}{XX}
  动量守恒                 & $m_1v_1+m_2v_2=m_1v_1'+m_2v_2'$                              \\
  恢复系数                 & $e=\frac{\text{分离相对速度}}{\text{接近相对速度}}$             \\
  一维恢复系数              & $v_2'-v_1'=e(v_1-v_2)$                                        \\
  约化质量                 & $\mu=\frac{m_1m_2}{m_1+m_2}$                                  \\
  碰撞损失机械能             & $\Delta E=\frac12\mu(1-e^2)(v_1-v_2)^2$                       \\
\end{tabularx}

\begin{formulanotebox}
碰撞问题首先看碰撞时间极短，外力冲量是否可以忽略。若外力冲量可以忽略，系统动量守恒；若碰撞过程光滑且没有能量损失，才进一步使用机械能守恒。\\
碰撞最推荐从相对速度分析，而不是直接背一长串碰撞后速度公式。恢复系数的本质是
\[
  \text{分离相对速度}=e\times\text{接近相对速度}.
\]
一维情况下，如果碰前$v_1>v_2$，两个物体正在接近，则
\[
  v_2'-v_1'=e(v_1-v_2).
\]
这个式子和动量守恒联立，已经足够解决绝大多数碰撞题。
\end{formulanotebox}

\begin{keypointbox}
碰撞按机械能损失分为三类：弹性碰撞、非弹性碰撞、完全非弹性碰撞。\\
弹性碰撞满足动量守恒和机械能守恒，恢复系数$e=1$。\\
非弹性碰撞满足动量守恒，但机械能减少，恢复系数$0<e<1$。\\
完全非弹性碰撞是机械能损失最大的碰撞，碰后两物体具有共同速度，恢复系数$e=0$。\\
约化质量
\[
  \mu=\frac{m_1m_2}{m_1+m_2}
\]
可以把两个物体的相对运动写得更紧凑。碰撞中真正改变的是相对运动，质心运动不受内力影响。
\end{keypointbox}

\begin{casetypebox}[title={\strut\textcolor{white}{\bfseries 常见题型：碰撞分类}}]
设两个物体沿同一直线碰撞，碰前速度为$v_1,v_2$，碰后速度为$v_1',v_2'$。只要碰撞时间内外力冲量可以忽略，就有
\[
  m_1v_1+m_2v_2=m_1v_1'+m_2v_2'.
\]
弹性碰撞还满足
\[
  \frac12m_1v_1^2+\frac12m_2v_2^2
  =\frac12m_1{v_1'}^2+\frac12m_2{v_2'}^2.
\]
但是实际解题更推荐把弹性碰撞写成相对速度反向且大小不变：
\[
  v_2'-v_1'=v_1-v_2.
\]
非弹性碰撞一般用恢复系数描述：
\[
  v_2'-v_1'=e(v_1-v_2),\qquad 0<e<1.
\]
完全非弹性碰撞中两物体碰后粘在一起，设共同速度为$v$：
\[
  m_1v_1+m_2v_2=(m_1+m_2)v.
\]
此时
\[
  v=\frac{m_1v_1+m_2v_2}{m_1+m_2}.
\]
这个共同速度其实就是系统质心速度。
\end{casetypebox}

\begin{casetypebox}[title={\strut\textcolor{white}{\bfseries 常见题型：恢复系数与约化质量}}]
恢复系数$e$只描述碰撞方向上的相对速度变化。以一维碰撞为例，碰前接近相对速度为$v_1-v_2$，碰后分离相对速度为$v_2'-v_1'$，所以
\[
  e=\frac{v_2'-v_1'}{v_1-v_2}.
\]
若$e=1$，相对速度只是反向，大小不变；若$e=0$，分离相对速度为0，两物体碰后不再分离。\\
用约化质量可以写出碰撞前相对运动的动能：
\[
  E_{\text{相}}=\frac12\mu(v_1-v_2)^2.
\]
碰后相对速度变为$e(v_1-v_2)$，相对运动动能变为
\[
  E_{\text{相}}'=\frac12\mu e^2(v_1-v_2)^2.
\]
因此碰撞损失的机械能为
\[
  \Delta E=\frac12\mu(1-e^2)(v_1-v_2)^2.
\]
完全非弹性碰撞时$e=0$，损失的正是全部相对运动动能：
\[
  \Delta E_{\max}=\frac12\mu(v_1-v_2)^2.
\]
这个结果说明：碰撞不会改变质心运动的动能，损失只来自相对运动。
\end{casetypebox}

\begin{casetypebox}[title={\strut\textcolor{white}{\bfseries 常见题型：一维碰撞的相对速度法}}]
一维碰撞不建议直接背两个碰后速度公式。更稳妥的步骤是：\\
第一，规定正方向，写动量守恒：
\[
  m_1v_1+m_2v_2=m_1v_1'+m_2v_2'.
\]
第二，写恢复系数：
\[
  v_2'-v_1'=e(v_1-v_2).
\]
第三，联立两个一次方程求$v_1',v_2'$。\\
如果题目给的是弹性碰撞，就令$e=1$；如果是完全非弹性碰撞，就令$e=0$，或者直接写$v_1'=v_2'$。\\
特殊情况下可以记结论：质量相等的一维弹性碰撞会交换速度。但这个结论只适用于一维正碰、质量相等、弹性碰撞，不要推广到二维或斜碰。
\end{casetypebox}

\begin{casetypebox}[title={\strut\textcolor{white}{\bfseries 常见题型：二维碰撞}}]
二维碰撞要先分解方向。光滑球碰撞时，碰撞冲量只沿两球球心连线方向，切向没有冲量。因此把速度分解为法向和切向：法向使用一维碰撞，切向速度保持不变。\\
设碰撞瞬间从物体1指向物体2的单位矢量为$\vec n$，碰前相对速度在法向上的分量为
\[
  u_n=(\vec v_1-\vec v_2)\cdot\vec n.
\]
若$u_n>0$，说明两物体正在沿法向接近。碰撞冲量大小为
\[
  J=(1+e)\mu u_n.
\]
于是碰后速度为
\[
  \vec v_1'=\vec v_1-\frac{J}{m_1}\vec n,\qquad
  \vec v_2'=\vec v_2+\frac{J}{m_2}\vec n.
\]
这等价于法向满足
\[
  (\vec v_2'-\vec v_1')\cdot\vec n
  =e(\vec v_1-\vec v_2)\cdot\vec n,
\]
而切向速度不变。\\
二维碰撞最容易错在把整个速度大小拿去套一维公式。真正发生突变的只有碰撞冲量方向上的速度分量，垂直于冲量方向的分量不参与正碰分析。
\end{casetypebox}

\begin{casetypebox}[title={\strut\textcolor{white}{\bfseries 常见题型：反冲运动}}]
反冲运动也是动量守恒和相对速度问题。系统内部有一部分物体向某个方向射出，剩余部分就会向相反方向运动。\\
设原系统初始静止，质量为$M$的主体和质量为$m$的喷出物分离，喷出物相对主体的速度大小为$u$。取喷出物方向为正，设主体速度为$V$，喷出物对地速度为$v$，则
\[
  M V+mv=0,
\]
\[
  v-V=u.
\]
联立得
\[
  v=\frac{M}{M+m}u,\qquad
  V=-\frac{m}{M+m}u.
\]
如果系统原来有共同初速度$v_0$，则可以先随质心一起平移，再叠加上面的相对运动结果。也就是说，总动量决定质心速度，相对速度决定两部分怎样分开。\\
反冲题不要一开始就猜哪个速度等于给定速度，题目给出的常常是“相对速度”。先判断速度是对地速度还是相对速度，再列动量守恒。
\end{casetypebox}

\begin{casetypebox}[title={\strut\textcolor{white}{\bfseries 常见题型：弹簧类碰撞}}]
弹簧类碰撞指两个物体之间通过轻弹簧相互作用，或者碰撞过程中可以等效为弹性形变储能。它和普通碰撞的相同点是：水平方向外力冲量为零时，系统动量守恒；不同点是：相互作用不是瞬间完成的，两个物体速度会连续变化。\\
设两物体质量分别为$m_1,m_2$，初速度为$v_1,v_2$，且$v_1>v_2$，后方物体追上前方物体并压缩弹簧。整个过程中
\[
  m_1v_1+m_2v_2=m_1v_1(t)+m_2v_2(t).
\]
如果弹簧和接触面都理想，没有能量损失，则机械能守恒：
\[
  \frac12m_1v_1^2+\frac12m_2v_2^2
  =\frac12m_1v_1(t)^2+\frac12m_2v_2(t)^2+E_p(t).
\]
弹簧压缩量最大时，两物体相对速度为零，因此
\[
  v_1=v_2=v_c.
\]
由动量守恒得共同速度
\[
  v_c=\frac{m_1v_1+m_2v_2}{m_1+m_2}.
\]
此时弹簧弹性势能最大：
\[
  E_{p\max}
  =\frac12m_1v_1^2+\frac12m_2v_2^2
  -\frac12(m_1+m_2)v_c^2
  =\frac12\mu(v_1-v_2)^2.
\]
弹簧重新恢复原长时，弹性势能又变回零。如果没有能量损失，最终效果等价于一次一维弹性碰撞：相对速度大小不变、方向反向：
\[
  v_2'-v_1'=v_1-v_2.
\]
\begin{center}
\begin{tikzpicture}[scale=0.95, >=Stealth]
  \draw[->] (0,0) -- (6.2,0) node[right] {$t$};
  \draw[->] (0,0) -- (0,3.0) node[above] {$v$};
  \draw[dashed] (2.5,0) -- (2.5,2.55);
  \node[below] at (2.5,0) {$t_c$};
  \node[above] at (2.5,2.55) {压缩最大，共速};

  \draw[thick, blue] (0.35,2.55) .. controls (1.3,2.45) and (1.8,1.75) .. (2.5,1.65)
    .. controls (3.3,1.55) and (4.0,0.95) .. (5.45,0.75);
  \draw[thick, red] (0.35,0.75) .. controls (1.3,0.85) and (1.8,1.55) .. (2.5,1.65)
    .. controls (3.3,1.75) and (4.0,2.35) .. (5.45,2.55);

  \draw[dotted] (0,1.65) -- (5.65,1.65);
  \node[left] at (0,1.65) {$v_c$};
  \node[blue, right] at (5.45,0.75) {$v_1(t)$};
  \node[red, right] at (5.45,2.55) {$v_2(t)$};
  \node[blue, left] at (0.35,2.55) {$v_1$};
  \node[red, left] at (0.35,0.75) {$v_2$};
  \node at (1.25,-0.35) {弹簧压缩};
  \node at (4.1,-0.35) {弹簧恢复};
\end{tikzpicture}
\end{center}
图像中，后方物体速度逐渐减小，前方物体速度逐渐增大；在弹簧压缩最大时两条速度曲线相交；之后弹簧恢复，两者继续交换能量。这个题型的关键是把“压缩最大”翻译成“相对速度为零”，把“恢复原长”翻译成“弹性势能为零”。
\end{casetypebox}
